
\documentclass{article}
\begin{document}

\subsection*{Estructura para resolver un juego de Stackelberg (dinámico)}

\subsubsection*{Paso 1. Definir el sistema dinámico}

\textbf{Estados:}
\[
x(t) \in \mathbb{R}^n
\]

\textbf{Dinámica:}
\[
\dot{x}(t) = f\big(x(t), u(t), v(t)\big), 
\qquad x(0) = x_0
\]

\subsubsection*{Paso 2. Definir los controles y conjuntos admisibles}

\[
u(t) \in U, \qquad v(t) \in V
\]

donde \(U, V\) son conjuntos cerrados (y convexos si se desea optimalidad clásica).

\subsubsection*{Paso 3. Definir los funcionales de costo}

\textbf{Líder:}
\[
J_L(u,v) = \int_0^T L_L\big(x(t),u(t),v(t)\big)\,dt
+ \Phi_L\big(x(T)\big)
\]

\textbf{Seguidor:}
\[
J_F(u,v) = \int_0^T L_F\big(x(t),u(t),v(t)\big)\,dt
+ \Phi_F\big(x(T)\big)
\]



Sí: \(v(t)\) es un \textbf{control adaptativo}.

En el modelo propuesto,
\[
v(t) = \text{control adaptativo del tumor metastásico}.
\]

El carácter adaptativo de \(v(t)\) se justifica por las siguientes razones:
\begin{itemize}
    \item responde a la presión terapéutica \(u(t)\);
    \item depende del estado tumoral \(P(t)\) y \(M(t)\);
    \item evoluciona dinámicamente en el tiempo;
    \item implica un costo interno de tipo metabólico y/o evolutivo.
\end{itemize}

En conjunto, estas características permiten interpretar a \(v(t)\) como un mecanismo de adaptación del sistema tumoral frente a la intervención terapéutica

\end{document}
